\documentclass[12pt]{article}
\usepackage[utf8]{inputenc}
\usepackage[T1]{fontenc}
\usepackage[margin=1in]{geometry}
\usepackage{setspace}
\usepackage{xcolor}
\usepackage[colorlinks=true, urlcolor=blue!60!black, linkcolor=blue!60!black]{hyperref}
\setlength{\parindent}{0em}
\setlength{\parskip}{0.45em}
\renewcommand{\baselinestretch}{1.0}
\begin{document}
\begin{center}
{\Large \textit{Plessy v.\ Ferguson}, 163 U.S.\ 537 (1896)\par}
\vspace{0.3em}
{\normalsize Case Brief --- One-Page Edition\par}
\end{center}

\textbf{Memory Jogger.} ``Separate but equal'' railway cars upheld over Harlan's color-blind dissent (\href{https://supreme.justia.com/cases/federal/us/163/537/}{Justia}).

\textbf{Facts.} Louisiana's Separate Car Act required ``equal but separate accommodations'' on intrastate railways. In a planned test case, Homer Plessy sat in a whites-only coach in 1892, refused to move, and was forcibly ejected and arrested.

\textbf{Procedural History.} Judge Ferguson rejected Plessy's constitutional plea; the Louisiana Supreme Court denied relief.

\textbf{Votes.} Affirmed 7--1: Brown for the Court; Harlan dissented; Brewer did not participate.

\textbf{Holding.} Mandated separate but equal intrastate accommodations violate neither the Thirteenth nor the Fourteenth Amendment, which secures political but not ``social'' equality, if reasonable under the police power.

\textbf{Reasoning.} A law drawing ``merely a legal distinction'' is not servitude, and any badge of inferiority arose ``solely'' from how Black citizens chose to read the law. Harlan: ``Every one knows'' the Act aimed to exclude Black passengers from white coaches; it was a ``badge of servitude.''

\textbf{Examples.} \emph{Same-Side:} segregated but nominally equal buses. \emph{Opposite-Side:} a whites-only law school with no equal alternative. \emph{Fence-Sitter:} segregation on interstate trains.

\textbf{Critique.} An anti-canon case, willfully blind to segregation's purpose of enforcing white supremacy; originalists still debate whether the framers forbade segregation, and scholars whether Harlan's dissent is color-blind or anti-caste.

\textbf{Key Quotation.} ``But in view of the constitution, in the eye of the law, there is in this country no superior, dominant, ruling class of citizens. There is no caste here. Our constitution is color-blind, and neither knows nor tolerates classes among citizens.'' --- Justice Harlan, dissenting, 163 U.S.\ at 559.

\end{document}
